% Personaliza \maketitle: 1+2 baselineskip, titulo en mayusculas, y las
% páginas de título internas (extratitle, uppertitleback, lowertitleback,
% dedication) definidas en 28-titlepages.tex.
%
% Orden de páginas:
%   extratitle (pág 1, bloque del 75% del ancho centrado con
%   \vspace*{7\baselineskip} antes)
%   → portada (impar) → titlebacks (reverso, twoside; el \vfill ancla el
%   lowertitleback al fondo, sin minipage) → dedication (página impar
%   siguiente, bloque del 50% del ancho alineado a la derecha como el dictum,
%   con \vspace*{7\baselineskip} antes).
% \titlepage@next/\titlepage@nextdouble (28-titlepages.tex): \clearpage +
% \thispagestyle{empty} (+ página impar en twoside), SIN el \setparsizes
% de \next@tpage (que dejaría \parindent a 0 en todo el documento).
%
% KOMA define \subtitle con \newcommand* (no-long): una línea en blanco en el
% argumento (párrafo del frontmatter |) rompe la compilación con 'Paragraph
% ended before \subtitle was complete'. Se redefine como long (\renewcommand
% sin estrella) para permitir subtítulos multilínea con párrafos.
% title-image (campo del frontmatter, solo LaTeX/PDF): imagen que sustituye al
% texto del título en la portada. graphicx NO venía cargado en el preamble.
\usepackage{graphicx}
\makeatletter
\renewcommand{\subtitle}[1]{\gdef\@subtitle{%
    #1%
}}%
\newcommand*{\@titleimage}{}%
\newcommand{\titleimage}[1]{\gdef\@titleimage{%
    #1%
}}%
\newcommand*{\@publishersimagelist}{}%
\newcommand{\publishersimage}[1]{%
  \ifx\@publishersimagelist\@empty
    \gdef\@publishersimagelist{#1}%
  \else
    \g@addto@macro\@publishersimagelist{,#1}%
  \fi
}%
\newcommand{\@renderpublishersimages}{%
  \@for\@pub@img:=\@publishersimagelist\do{%
    \titleimagerender[150pt]{\@pub@img}%
  }%
}%
\newcommand{\@renderpublishersimages@small}{%
  \@for\@pub@img:=\@publishersimagelist\do{%
    \titleimagerender[0.25\textwidth]{\@pub@img}%
  }%
}%

% Render de la imagen de portada con ancho MÁXIMO del 80% del textwidth: se
% mide el ancho natural con un \sbox; si supera el 80% se escala conservando
% la proporción, si no se muestra a tamaño natural (no se amplía).
\newsavebox{\titleimagebox}
% Render de la imagen de portada con ancho MÁXIMO configurable: #1 (opcional,
% por defecto 0.8\textwidth) es el ancho máximo y #2 la ruta. Se mide el ancho
% natural con un \sbox; si supera el máximo se escala conservando la
% proporción, si no se muestra a tamaño natural (no se amplía). La portada usa
% el default; la página de extratitle pasa [\extratitlewidth] (100% del ancho
% del bloque).
\newcommand{\titleimagerender}[2][0.8\textwidth]{%
  \sbox{\titleimagebox}{\includegraphics{#2}}%
  \ifdim\wd\titleimagebox>#1
    \includegraphics[width=#1,keepaspectratio]{#2}%
  \else
    \usebox{\titleimagebox}%
  \fi
}

% Bloques de las páginas de título con la geometría del dictum (21-dictum.tex):
% list con márgenes dados (#1 izquierdo, #2 derecho) y sin espaciado extra.
% #3 es el estilo del párrafo (\centering, \raggedright o vacío = justificado):
% el extratitle centra su texto, la dedication lo justifica.
% Geometría propia para no crear dependencia con 21-dictum.tex.
\newcommand*{\extratitlewidth}{0.75\textwidth}%
\newcommand*{\dedicationwidth}{0.5\textwidth}%
\newcommand*{\colophonwidth}{0.75\textwidth}%
\newcommand{\titlepageblock}[4]{%
  \par
  \begin{list}{}{%
      \leftmargin=#1
      \rightmargin=#2
      \itemsep=\z@
      \parsep=\z@
      \partopsep=\z@
      \topsep=\z@
      \listparindent=\z@
      \itemindent=\z@
    }%
    \item\relax
    {#3 #4\par}%
  \end{list}%
}

% Dos páginas en blanco (1 y 2) antes de la página de extratitle: el libro
% abre con hojas de guarda y la extratitle queda en la página 3 (impar).
% Cada \null evita que la página sea vacía (\clearpage de una página sin
% contenido no la emite en algunos casos). Sin número: los layers de
% scrlayer-scrpage están vacíos hasta que el comando de página se define.
% \titlepageguards: true si las hojas de guarda existen — 30-endpapers.tex
% la consulta para dibujar el endpaper SOLO sobre la página 1 cuando esta es
% la hoja en blanco (sin guardas, sin endpaper en ningún caso).
\newif\iftitlepageguards
\newcommand{\titlepageblanks}{%
  \titlepageguardstrue
  \null\clearpage
  \null\clearpage
}
\renewcommand{\maketitle}{%
  % Página de extratitle (pág 3, centrada horizontal y verticalmente): existe
  % si hay extratitle, title-image o frontispiece. Con solo frontispiece, la
  % página se llena por defecto con el título.
  % Centrado vertical con \vbox to \textheight: el \vspace*{\fill} entre
  % bloques dejaba el contenido ~9-17pt fuera del centro (interacción con el
  % page builder y el strut); dentro del vbox los \vfill reparten el espacio
  % exactamente. No confina el ancho (no es minipage).
  \ifx\@extratitle\@empty
    \ifx\@frontispiece\@empty
      % sin página de extratitle
    \else
      \titlepageblanks
      \vbox to \textheight{%
        \vfill
        {\centering\parindent\z@\@title\par}%
        \vfill
      }%
      \titlepage@next
    \fi
  \else
    \titlepageblanks
    \vbox to \textheight{%
      \vspace*{10\baselineskip}%
      \titlepageblock
        {\dimexpr(\linewidth-\extratitlewidth)/2\relax}%
        {\dimexpr(\linewidth-\extratitlewidth)/2\relax}%
        {\centering}%
        {\@extratitle}%
      \vfill
    }%
    \ifx\@frontispiece\@empty
      \titlepage@nextdouble
    \else
      \titlepage@next
    \fi
  \fi
  % Frontispicio: página anterior a la portada, con relleno vertical antes
  % del contenido para que quede anclado al fondo.
  \ifx\@frontispiece\@empty\else
    \vspace*{\fill}%
    {\centering\parindent\z@\@frontispiece\par}%
    \titlepage@next
  \fi
  \thispagestyle{empty}%
  \vspace*{\baselineskip}%
  \ifx\@author\@empty
    \vskip 2\baselineskip%
  \else
    {\centering\usekomafont{author}{\renewcommand{\and}{\unskip, \ignorespaces}\@author\par}}%
    \vskip 1\baselineskip
  \fi
  \ifx\@titleimage\@empty
    {\centering\usekomafont{title}{\MakeUppercase{\@title}\par}}%
  \else
    {\centering\titleimagerender{\@titleimage}\par}%
  \fi
  \ifx\@subtitle\@empty\else
    \vskip 1\baselineskip
    {\centering\usekomafont{subtitle}\parindent\z@\@subtitle\par}%
  \fi
  \ifx\@subject\@empty\else
    \vskip 1\baselineskip
    {\centering\parindent\z@\@subject\par}%
  \fi
  \ifx\@date\@empty\else
    \vskip 1\baselineskip
    {\centering\usekomafont{date}{\@date\par}}%
  \fi
  \ifx\@titlehead\@empty\else
    \vskip 2\baselineskip
    {\centering\parindent\z@\@titlehead\par}%
  \fi
  \ifx\@publishersimagelist\@empty
    \ifx\@publishers\@empty\else
      \vskip 2\baselineskip
      {\centering\parindent\z@\@publishers\par}%
    \fi
  \else
    % publisher-image definido: la imagen sustituye al texto de publishers
    % (máx. 150pt ≈ 150px a 72dpi, ajustable). Soporta múltiples imágenes.
    \vskip 2\baselineskip
    {\centering\@renderpublishersimages\par}%
  \fi
  \if@twoside
    \@tempswatrue
    \ifx\@uppertitleback\@empty\ifx\@lowertitleback\@empty
      \@tempswafalse
    \fi\fi
    \if@tempswa
      \titlepage@next
      {\titlebackformat\@uppertitleback\par}%
      \vfill
      {\titlebackformat\@lowertitleback\par}%
    \fi
  \fi
  \ifx\@dedication\@empty\else
    \titlepage@nextdouble
    \vspace*{7\baselineskip}%
    \titlepageblock
      {\dimexpr\linewidth-\dedicationwidth\relax}%
      {\z@}%
      {}%
      {\@dedication}%
    \clearpage
  \fi
}
\makeatother
